\section{Motivation}

\subsection{Motivating Example}

\autoref{tab:example} lists the requirements of a formal system specifying an algorithm that searches for the minimum value of $f(j)$ over all valid $j$ in the domain of the function $f$. 
In this example, the requirements are divided into two categories: $EQP$ (equipment) and $FUN$ (function). 
The equipment requirements define the constants, sets, and axioms available to the system, while the function requirements specify its intended behavior. 
In this example, $EQP$--1 and $EQP$--2 introduce the constants $n$ and $f$.
$FUN$--1 states that upon termination, the algorithm must return an index $j \in dom(f)$ such that $f(j)$ is the minimum value in the range of $f$.
Requirements $FUN$--2 through $FUN$--6 describe the iterative searching process to achieve this goal.

\begin{table}[h!]
\centering
\footnotesize
\begin{tabular}{l p{0.85\linewidth}}
\toprule
\textbf{ID} & \textbf{Description} \\
\midrule
EQP-1 & Provides a finite natural number function $f$ defined on the domain $\{0..n-1\}$. \\
EQP-2 & Provides $n$ as the size of the domain of $f$. \\
FUN-1 & The final condition requires $j \in \text{dom}(f)$ and $f(j) = \min(\text{ran}(f))$. \\
FUN-2 & The system has a boolean variable \texttt{searching} to indicate whether the minimum has been fully identified. Two indices are maintained: $i$ as the scanning index and $j$ as the current candidate for the minimum. \\
FUN-3 & Initialization sets i = 1 and j = 0, with searching = true. \\
FUN-4 & While \texttt{searching} is true, if $i \neq n$ and $f(i) < f(j)$, update $j := i$ and increment $i$. \\
FUN-5 & While \texttt{searching} is true, if $i \neq n$ and $f(i) \geq f(j)$, keep $j$ unchanged and increment $i$. \\
FUN-6 & When $i = n$, \texttt{searching} becomes false, marking the completion of searching. \\
\bottomrule
\end{tabular}
\caption{Equipment (EQP) and Function (FUN) requirements of minimum-searching model.}
\label{tab:example}
\end{table}


\begin{figure}[t]
\includegraphics[width=\linewidth]{figures/example_v4.drawio.pdf}
\caption{Simplified Event-B models for the motivating example generated by GPT-5, Cursor, and PAT-Agent. }
\label{fig:example_new}
\centering
\end{figure}

\begin{figure}[t]
\includegraphics[width=0.85\linewidth]{figures/running_example_simplified.drawio.pdf}
\caption{The abstract and concrete formal models constructed by Event-B Agent for the motivating example.}
\label{fig:running_example}
\centering
\end{figure}

As discussed in Section~\ref{sec:intro}, we select a set of existing LLM-based autoformalization approaches, namely GPT-5 with medium reasoning, Cursor, and PAT-Agent, and adapt them for the task of formal model construction.
Their performances are illustrated in \autoref{fig:example_new}, while the formal models constructed by Event-B Agent are visualized in \autoref{fig:running_example}.
The ``Context'' block in \autoref{fig:example_new} represents the shared context across the models generated by various methods, corresponding to the EQP requirements.
The table on the left of the figure shows how each method formalizes requirement $FUN$--1 to express the algorithm termination condition, well-definedness, and correctness condition.
In model $M_{Cursor}$ generated by Cursor, the $INITIALISATION$ event sets $j:=0$ and $searching:=FALSE$, violating $FUN$--1 unless $f(0)=min(ran(f))$.
In contrast, the remaining baselines do not mistakenly impose termination constraints during initialization, but exhibit other limitations.
GPT-5 produces model $M_{GPT-5}$ with $FUN$--1: $searching=FALSE \; \land \; j \in dom(f) \Rightarrow f(j)=min(ran(f))$, where the termination condition is not required to hold in the initial state.
However, in the $stop$ event, the guards do not restrict $j$ or $f(j)$, allowing the algorithm to terminate with an arbitrary $j$ and violate $FUN$--1.
PAT-Agent, a framework specialized for formal model synthesis, produces a stronger model $M_{PAT-Agent}$ in which the guards of $stop$ enforce $valAtJ=minRanF$, thereby satisfying $FUN$--1.
Nonetheless, this model contains subtle flaws.
In $INITIALISATION$, both $valAtI$ and $valAtJ$ are set to $3$, and subsequent events only update $valAtJ:=valAtI$ without modifying $valAtI$. 
As a result, the guard of $stop$ implies that the constant $minRanF$ must always be $3$. 
From the perspective of model checking, this system is deemed correct, but the correctness does not generalize beyond the single case $minRanF=3$.

The models generated by Event-B Agent, $M_{abstract}$ and $M_{concrete}$, resolve this issue through refinement.
The models are shown in \autoref{fig:running_example}, where events are represented by state transitions.
We use $[G]|A$ to denote an event with guards $G$ and actions $A$.
Since $FUN$--1 depends only on $j$ and $f$, the abstract model omits auxiliary variables such as $i$ and $searching$, and focuses solely on proving $FUN$--1.
$FUN$--1 is vacuously true for $INITIALISATION$ and $stop\_empty$ events.
For the $stop$ event, $FUN$--1 is preserved because the event’s assignments ensure $has\_j = TRUE \land j \in dom(f)$ in the post-state, and the guard enforces $f(j) = min(ran(f))$.
By construction of the refinement, the refined events $INITIALISATION$, $stop$, and $stop\_empty$, strengthen the guards of their abstract counterparts and simulate the abstract assignments, thereby preserving the invariant in $M_{\text{concrete}}$.

This example highlights three key insights motivating the design of Event-B Agent.
First, refinement reduces modeling and proving complexity by abstracting away irrelevant details. 
Second, when models are inaccurate, a systematic repair framework is required, e.g., updating the actions of $INITIALISATION$ in $M_{Cursor}$ and the guards of $stop$ in $M_{GPT-5}$.
Finally, model checking alone is insufficient to guarantee correctness, and proof information must be exploited in model synthesis and repair.
Accordingly, we propose Event-B Agent, which addresses the formal model synthesis and repair problem by incorporating refinement, a novel repair component, and proof information.


\subsection{Problem Statement} \label{sec:problem_statement}

Let $M$ denote the set of formulas extracted from a formal model, $\pi$ denote a proof artifact witnessing that a proof obligation (PO) $\varphi$ is derivable from $M$ under a sound proof system $\mathcal{R}$.
Then $M$ is said to discharge $\varphi$ if and only if there exists $\pi$ that witnesses $M \vdash_{\mathcal{R}} \varphi$.
We abstract formal development as a process over model construction and proof derivation, where both may be updated independently or simultaneously within each step.
Formally,
\[
(M^0, \pi_0) \rightsquigarrow (M^1, \pi_1) \rightsquigarrow ... \rightsquigarrow (M^n, \pi_n)
\text{ s.t. }
\forall \; t \in \{0,\dots,n\}, \pi_t \text{ is a valid proof for } M^t \vdash_{\mathcal{R}} \varphi_t,
\]
where $M^0 = \varnothing$, $\pi_0 = \epsilon$, $n$ is the number of POs, and the step $(M^0, \pi_0) \rightsquigarrow (M^1, \pi_1)$ synthesizes the initial model $M^1$ given natural language requirements and discharges $\pi_1$.
Each pair $(M^t, \pi_t)$ represents an intermediate model $M^t$ with a proof $\pi_t$ discharged at step $t$, allowing updates to both $M^t$ and $\pi_t$, while ensuring soundness at each step through $\pi_t$ under the sound proof system $\mathcal{R}$.

We position our work relative to two lines of research.
LLM proof assistants assume a fixed model $M$ and search for valid $\pi$, but cannot construct $M^0 \rightsquigarrow M^1$ from natural language requirements or revise $M^t$ during development.
In contrast, LLM-based autoformalization methods iteratively construct models $M^0 \rightsquigarrow \cdots \rightsquigarrow M^k$ without proof-based verification or repair, relying primarily on bounded model checking with limited correctness guarantees.
Neither line of work formulates formal development as a coordinated process of model construction and proof derivation.

Furthermore, refinement should be supported to incrementally develop the models.
Let $M_i^t$ denote the $t$-th model at refinement level $i$, where $i=\{1,\dots,m\}$.
For each $i < m$, the last model at level $i$ is refined by the first model at level $i+1$, with refinement soundness ensured by POs.
Eventually, $M_m^n$ is produced as the final model.
Existing LLM-based autoformalization does not support refinement in this sense, and LLM proof assistants operate over fixed models without refinement.